% !TeX program = XeLaTeX
% !TeX root = ../vedamantrabook.tex
\chapt{जयादि-होमः}
\label{sec:start_Jayadi}

\newcommand{\idannamama}[1]{\uwave{#1}~इदं न मम॥}
\newcommand{\namama}[1]{\uwave{#1}~न मम॥}

एतत्कर्मसमृद्ध्यर्थं जयादिहोमं करिष्ये॥

१. चि॒त्तं च॒ स्वाहा᳚।
\namama{चित्तायेदं}

२. चित्ति॑श्च॒ स्वाहा᳚।
\namama{चित्त्या इदं}

३. आकू॑तं च॒ स्वाहा᳚।
\namama{आकूतायेदं}

४. आकू॑तिश्च॒ स्वाहा᳚।
\namama{आकूत्या इदं}

५. विज्ञा॑तं च॒ स्वाहा᳚।
\namama{विज्ञातायेदं}

६. वि॒ज्ञानं॑ च॒ स्वाहा᳚।
\namama{विज्ञानायेदं}

७. मन॑श्च॒ स्वाहा᳚।
\namama{मनस इदं}

८. शक्व॑रीश्च॒ स्वाहा᳚।
\namama{शक्वरीभ्य इदं}

९. दर्श॑श्च॒ स्वाहा᳚।
\namama{दर्शायेदं}

१०. पू॒र्णमा॑सश्च॒ स्वाहा᳚।
\namama{पूर्णमासायेदं}

११. बृ॒हच्च॒ स्वाहा᳚।
\namama{बृहत इदं}

१२. र॒थ॒न्त॒रं च॒ स्वाहा᳚।
\namama{रथन्तरायेदं}

१३. प्र॒जाप॑ति॒र्जया॒निन्द्रा॑य॒ वृष्णे॒ प्राय॑च्छदु॒ग्रः पृ॑त॒नाज्ये॑षु॒ तस्मै॒ विशः॒ सम॑नमन्त॒ सर्वाः॒ स उ॒ग्रः स हि॒ हव्यो॑ ब॒भूव॒ स्वाहा᳚।
\namama{प्रजापतय इदं}


\clearpage

\dnsub{अभ्यातानाः}
\begingroup
\setlength{\ULdepth}{1ex}

१. \uwave{अ॒ग्निर्भू॒ताना॒मधि॑पतिः॒} स मा॑ऽवत्व॒स्मिन् ब्रह्म॑न्न॒स्मिन् क्ष॒त्रे᳚ऽस्यामा॒शिष्य॒स्यां पु॑रो॒धाया॑म॒स्मिन् कर्म॑न्न॒स्यां दे॒वहू᳚त्या॒ꣴ॒ स्वाहा᳚।
\namama{अग्नय इदं}

२. \uwave{इन्द्रो᳚ ज्ये॒ष्ठाना॒मधि॑पतिः॒} स मा॑ऽवत्व॒स्मिन् ब्रह्म॑न्न॒स्मिन् क्ष॒त्रे᳚ऽस्यामा॒शिष्य॒स्यां पु॑रो॒धाया॑म॒स्मिन् कर्म॑न्न॒स्यां दे॒वहू᳚त्या॒ꣴ॒ स्वाहा᳚।
\namama{इन्द्रायेदं}

३. \uwave{य॒मः पृ॑थि॒व्या अधि॑पतिः॒} स मा॑ऽवत्व॒स्मिन् ब्रह्म॑न्न॒स्मिन् क्ष॒त्रे᳚ऽस्यामा॒शिष्य॒स्यां पु॑रो॒धाया॑म॒स्मिन् कर्म॑न्न॒स्यां दे॒वहू᳚त्या॒ꣴ॒ स्वाहा᳚।
\namama{यमायेदं}

४. \uwave{वा॒युर॒न्तरि॑क्ष॒स्याधि॑पतिः॒} स मा॑ऽवत्व॒स्मिन् ब्रह्म॑न्न॒स्मिन् क्ष॒त्रे᳚ऽस्यामा॒शिष्य॒स्यां पु॑रो॒धाया॑म॒स्मिन् कर्म॑न्न॒स्यां दे॒वहू᳚त्या॒ꣴ॒ स्वाहा᳚।
\namama{वायव इदं}

५. \uwave{सूर्यो॑ दि॒वोऽधि॑पतिः॒} स मा॑ऽवत्व॒स्मिन् ब्रह्म॑न्न॒स्मिन् क्ष॒त्रे᳚ऽस्यामा॒शिष्य॒स्यां पु॑रो॒धाया॑म॒स्मिन् कर्म॑न्न॒स्यां दे॒वहू᳚त्या॒ꣴ॒ स्वाहा᳚।
\namama{सूर्यायेदं}

६. \uwave{च॒न्द्रमा॒ नक्ष॑त्राणा॒मधि॑पतिः॒} स मा॑ऽवत्व॒स्मिन् ब्रह्म॑न्न॒स्मिन् क्ष॒त्रे᳚ऽस्यामा॒शिष्य॒स्यां पु॑रो॒धाया॑म॒स्मिन् कर्म॑न्न॒स्यां दे॒वहू᳚त्या॒ꣴ॒ स्वाहा᳚।
\namama{चन्द्रमस इदं}

७. \uwave{बृह॒स्पति॒र्ब्रह्म॒णोऽधि॑पतिः॒} स मा॑ऽवत्व॒स्मिन् ब्रह्म॑न्न॒स्मिन् क्ष॒त्रे᳚ऽस्यामा॒शिष्य॒स्यां पु॑रो॒धाया॑म॒स्मिन् कर्म॑न्न॒स्यां दे॒वहू᳚त्या॒ꣴ॒ स्वाहा᳚।
\namama{बृहस्पतय इदं}

८. \uwave{मि॒त्रः स॒त्याना॒मधि॑पतिः॒} स मा॑ऽवत्व॒स्मिन् ब्रह्म॑न्न॒स्मिन् क्ष॒त्रे᳚ऽस्यामा॒शिष्य॒स्यां पु॑रो॒धाया॑म॒स्मिन् कर्म॑न्न॒स्यां दे॒वहू᳚त्या॒ꣴ॒ स्वाहा᳚।
\namama{मित्रायेदं}

९. \uwave{वरु॑णो॒ऽपामधि॑पतिः॒} स मा॑ऽवत्व॒स्मिन् ब्रह्म॑न्न॒स्मिन् क्ष॒त्रे᳚ऽस्यामा॒शिष्य॒स्यां पु॑रो॒धाया॑म॒स्मिन् कर्म॑न्न॒स्यां दे॒वहू᳚त्या॒ꣴ॒ स्वाहा᳚।
\namama{वरुणायेदं}

१०. \uwave{स॒मु॒द्रः स्रो॒त्याना॒मधि॑पतिः॒} स मा॑ऽवत्व॒स्मिन् ब्रह्म॑न्न॒स्मिन् क्ष॒त्रे᳚ऽस्यामा॒शिष्य॒स्यां पु॑रो॒धाया॑म॒स्मिन् कर्म॑न्न॒स्यां दे॒वहू᳚त्या॒ꣴ॒ स्वाहा᳚।
\namama{समुद्रायेदं}

११. \uwave{अन्न॒ꣳ॒ साम्रा᳚ज्याना॒मधि॑पति॒} तन्मा॑वत्व॒स्मिन् ब्रह्म॑न्न॒स्मिन् क्ष॒त्रे᳚ऽस्यामा॒शिष्य॒स्यां पु॑रो॒धाया॑म॒स्मिन् कर्म॑न्न॒स्यां दे॒वहू᳚त्या॒ꣴ॒ स्वाहा᳚।
\namama{अन्नायेदं}

१२. \uwave{सोम॒ ओष॑धीना॒मधि॑पतिः॒} स मा॑ऽवत्व॒स्मिन् ब्रह्म॑न्न॒स्मिन् क्ष॒त्रे᳚ऽस्यामा॒शिष्य॒स्यां पु॑रो॒धाया॑म॒स्मिन् कर्म॑न्न॒स्यां दे॒वहू᳚त्या॒ꣴ॒ स्वाहा᳚।
\namama{सोमायेदं}

१३. \uwave{स॒वि॒ता प्र॑स॒वाना॒मधि॑पतिः॒} स मा॑ऽवत्व॒स्मिन् ब्रह्म॑न्न॒स्मिन् क्ष॒त्रे᳚ऽस्यामा॒शिष्य॒स्यां पु॑रो॒धाया॑म॒स्मिन् कर्म॑न्न॒स्यां दे॒वहू᳚त्या॒ꣴ॒ स्वाहा᳚।
\namama{सवित्र इदं}

१४. \uwave{रु॒द्रः प॑शू॒नामधि॑पतिः॒} स मा॑ऽवत्व॒स्मिन् ब्रह्म॑न्न॒स्मिन् क्ष॒त्रे᳚ऽस्यामा॒शिष्य॒स्यां पु॑रो॒धाया॑म॒स्मिन् कर्म॑न्न॒स्यां दे॒वहू᳚त्या॒ꣴ॒ स्वाहा᳚।
\namama{रुद्रायेदं}
[अप उपस्पृश्य]

१५. \uwave{त्वष्टा॑ रू॒पाणा॒मधि॑पतिः॒} स मा॑ऽवत्व॒स्मिन् ब्रह्म॑न्न॒स्मिन् क्ष॒त्रे᳚ऽस्यामा॒शिष्य॒स्यां पु॑रो॒धाया॑म॒स्मिन् कर्म॑न्न॒स्यां दे॒वहू᳚त्या॒ꣴ॒ स्वाहा᳚।
\namama{त्वष्ट्र इदं}

१६. \uwave{विष्णुः॒ पर्व॑ताना॒मधि॑पतिः॒} स मा॑ऽवत्व॒स्मिन् ब्रह्म॑न्न॒स्मिन् क्ष॒त्रे᳚ऽस्यामा॒शिष्य॒स्यां पु॑रो॒धाया॑म॒स्मिन् कर्म॑न्न॒स्यां दे॒वहू᳚त्या॒ꣴ॒ स्वाहा᳚।
\namama{विष्णव इदं}

१७. \uwave{म॒रुतो॑ ग॒णाना॒मधि॑पतय॒स्ते} मा॑ऽवन्त्व॒स्मिन् ब्रह्म॑न्न॒स्मिन् क्ष॒त्रे᳚ऽस्यामा॒शिष्य॒स्यां पु॑रो॒धाया॑म॒स्मिन् कर्म॑न्न॒स्यां दे॒वहू᳚त्या॒ꣴ॒ स्वाहा᳚।
\namama{मरुद्भ्य इदं}

१८. \uwave{पित॑रः पितामहाः परेऽवरे॒ तता॑॑स्ततामहा इ॒ह मा॑ऽवत।} अ॒स्मिन् ब्रह्म॑न्न॒स्मिन् क्ष॒त्रे᳚ऽस्यामा॒शिष्य॒स्यां पु॑रो॒धाया॑म॒स्मिन् कर्म॑न्न॒स्यां दे॒वहू᳚त्या॒ꣴ॒ स्वाहा᳚।
\namama{पितृभ्य इदं}
[अप उपस्पृश्य]
\endgroup

\dnsub{राष्ट्रभृतः}
\begingroup
\setlength{\ULdepth}{1ex}

१. ऋ॒ता॒षाडृ॒तधा॑मा॒\uwave{ऽग्नि}र्ग॑न्ध॒र्वस्त\uwave{स्यौष॑धयो}ऽफ्स॒रस॒ \uwave{ऊर्जो॒} नाम॒
स इ॒दं ब्रह्म॑ क्ष॒त्रं पा॑तु॒ ता इ॒दं ब्रह्म॑ क्ष॒त्रं पा᳚न्तु॒ तस्मै॒ स्वाहा᳚।\\
\namama{अग्नये गन्धर्वायेदं}\\
ताभ्यः॒ स्वाहा᳚।
\namama{ओषधीभ्य अफ्सरोभ्य इदं}

२. स॒ꣳ॒हि॒तो वि॒श्वसा॑मा॒ \uwave{सूर्यो॑} गन्ध॒र्वस्तस्य॒ \uwave{मरी॑चयो}ऽफ्स॒रस॑ \uwave{आ॒युवो॒} नाम॒
स इ॒दं ब्रह्म॑ क्ष॒त्रं पा॑तु॒ ता इ॒दं ब्रह्म॑ क्ष॒त्रं पा᳚न्तु॒ तस्मै॒ स्वाहा᳚।\\
\namama{सूर्याय गन्धर्वायेदं}\\
ताभ्यः॒ स्वाहा᳚।
\namama{मरीचिभ्योऽफ्सरोभ्य इदं}

३. सु॒षु॒म्नः सूर्य॑रश्मि\uwave{श्च॒न्द्रमा॑} गन्ध॒र्वस्तस्य॒ \uwave{नक्ष॑त्राण्य}फ्स॒रसो॑ \uwave{बे॒कुर॑यो॒} नाम॒
स इ॒दं ब्रह्म॑ क्ष॒त्रं पा॑तु॒ ता इ॒दं ब्रह्म॑ क्ष॒त्रं पा᳚न्तु॒ तस्मै॒ स्वाहा᳚।\\
\namama{चन्द्रमसे गन्धर्वायेदं}\\
ताभ्यः॒ स्वाहा᳚।
\namama{नक्षत्रेभ्योऽफ्सरोभ्य इदं}

४. भु॒ज्युः सु॑प॒र्णो \uwave{य॒ज्ञो} ग॑न्ध॒र्वस्तस्य॒ \uwave{दक्षि॑णा} अफ्स॒रसः॑ \uwave{स्त॒वा} नाम॒
स इ॒दं ब्रह्म॑ क्ष॒त्रं पा॑तु॒ ता इ॒दं ब्रह्म॑ क्ष॒त्रं पा᳚न्तु॒ तस्मै॒ स्वाहा᳚।\\
\namama{यज्ञाय गन्धर्वायेदं}\\
ताभ्यः॒ स्वाहा᳚।
\namama{दक्षिणाभ्योऽफ्सरोभ्य इदं}

५. प्र॒जाप॑तिर्वि॒श्वक॑र्मा॒ \uwave{मनो॑} गन्ध॒र्वस्तस्य॑\uwave{र्ख्सा॒मान्य॑}फ्स॒रसो॒ \uwave{वह्न॑यो॒} नाम॒
स इ॒दं ब्रह्म॑ क्ष॒त्रं पा॑तु॒ ता इ॒दं ब्रह्म॑ क्ष॒त्रं पा᳚न्तु॒ तस्मै॒ स्वाहा᳚।\\
\namama{मनसे गन्धर्वायेदं}\\
ताभ्यः॒ स्वाहा᳚।
\namama{ऋख्सामेभ्योऽफ्सरोभ्य इदं}

६. इ॒षि॒रो वि॒श्वव्य॑चा॒ \uwave{वातो॑} गन्ध॒र्वस्तस्या\uwave{ऽऽपो᳚}ऽफ्स॒रसो॑ \uwave{मु॒दा} नाम॒
स इ॒दं ब्रह्म॑ क्ष॒त्रं पा॑तु॒ ता इ॒दं ब्रह्म॑ क्ष॒त्रं पा᳚न्तु॒ तस्मै॒ स्वाहा᳚।\\
\namama{वाताय गन्धर्वायेदं}\\
ताभ्यः॒ स्वाहा᳚।
\namama{अद्भ्योऽफ्सरोभ्य इदं}

७. \uwave{भुव॑नस्य पते॒} यस्य॑ त उ॒परि॑ गृ॒हा इ॒ह च॑।
स नो॑ रा॒स्वाज्या॑निꣳ रा॒यस्पोषꣳ॑ सु॒वीर्यꣳ॑ संवत्स॒रीणाꣴ॑ स्व॒स्तिꣴ स्वाहा᳚॥
\namama{भुवनस्य पत्य इदं}

८. प॒र॒मे॒ष्ठ्यधि॑पति\uwave{र्मृ॒त्यु}र्ग॑न्ध॒र्वस्तस्य॒ \uwave{विश्व॑म}फ्स॒रसो॒ \uwave{भुवो॒} नाम॒
स इ॒दं ब्रह्म॑ क्ष॒त्रं पा॑तु॒ ता इ॒दं ब्रह्म॑ क्ष॒त्रं पा᳚न्तु॒ तस्मै॒ स्वाहा᳚।\\
\namama{मृत्यवे गन्धर्वायेदं}\\
ताभ्यः॒ स्वाहा᳚।
\namama{विश्वस्मा अफ्सरोभ्य इदं}

९. सु॒क्षि॒तिः सुभू॑तिर्भद्र॒कृत् सुव॑र्वान् \uwave{प॒र्जन्यो॑} गन्ध॒र्वस्तस्य॑ \uwave{वि॒द्युतो॑॑}ऽफ्स॒रसो॒ \uwave{रुचो॒} नाम॒
स इ॒दं ब्रह्म॑ क्ष॒त्रं पा॑तु॒ ता इ॒दं ब्रह्म॑ क्ष॒त्रं पा᳚न्तु॒ तस्मै॒ स्वाहा᳚।
\namama{पर्जन्याय गन्धर्वाय इदं}\\
ताभ्यः॒ स्वाहा᳚।
\namama{विद्युद्भ्योऽफ्सरोभ्य इदं}

१०. दू॒रे हे॑तिरमृड॒यो \uwave{मृ॒त्यु}र्ग॑न्ध॒र्वस्तस्य॑ \uwave{प्र॒जा} अ॑फ्स॒रसो॑ \uwave{भी॒रुवो॒} नाम॒
स इ॒दं ब्रह्म॑ क्ष॒त्रं पा॑तु॒ ता इ॒दं ब्रह्म॑ क्ष॒त्रं पा᳚न्तु॒ तस्मै॒ स्वाहा᳚।\\
\namama{मृत्यवे गन्धर्वायेदं}\\
ताभ्यः॒ स्वाहा᳚।
\namama{प्रजाभ्योऽफ्सरोभ्य इदं}

११. चारुः॑ कृपणका॒शी \uwave{कामो॑} गन्ध॒र्वस्तस्या॒\-\uwave{ऽऽधयो᳚}ऽफ्स॒रसः॑ \uwave{शो॒चय॑न्ती॒र्}नाम॒
स इ॒दं ब्रह्म॑ क्ष॒त्रं पा॑तु॒ ता इ॒दं ब्रह्म॑ क्ष॒त्रं पा᳚न्तु॒ तस्मै॒ स्वाहा᳚।\\
\namama{कामाय गन्धर्वायेदं}\\
ताभ्यः॒ स्वाहा᳚।
\namama{आधिभ्योऽफ्सरोभ्य इदं}

१२. स नो॑ भुवनस्य पते॒ यस्य॑ त उ॒परि॑ गृ॒हा इ॒ह च॑।
उ॒रु ब्रह्म॑णे॒ऽस्मै क्ष॒त्राय॒ महि॒ शर्म॑ यच्छ॒ स्वाहा᳚॥\\
\namama{भुवनस्य पत्य ब्रह्मण इदं}
\endgroup

\dnsub{प्रजापत्या}
[उपांशुना]

प्रजा॑पते॒ न त्वदे॒तान्य॒न्यो विश्वा॑ जा॒तानि॒ परि॒ ता ब॑भू॒व।

यत्का॑मास्ते जुहु॒मस्तन्नो॑ अस्तु व॒यꣴ स्या॑म॒ पत॑यो रयी॒णाꣴ स्वाहा᳚।

प्रजापतय इदं न मम॥

\dnsub{व्याहृतयः}
भूः स्वाहा᳚। अग्नय इदं न मम॥

भुवः॒ स्वाहा᳚। वायव इदं न मम॥

सुवः॒ स्वाहा᳚। सूर्याय इदं न मम॥

\dnsub{सौविष्टकृती}
यद॑स्य॒ कर्म॒णोऽत्यरी॑रिचं॒ यद्वा॒ न्यू॑नमि॒हाक॑रम्।

अ॒ग्नि॒ष्टत् स्वि॑ष्ट॒कृद् वि॒द्वान् सर्व॒ꣴ॒ स्वि॑ष्ट॒ꣳ॒ सुहु॑तं करोतु॒ स्वाहा᳚।

अग्नये स्विष्टकृत इदं न मम॥

\hyperref[sec:start_Jayadi]{\closesection}